% 00-resumo.tex | Abstract estruturado da Production: ate 250 palavras, cinco partes rotuladas
% (Paper aims; Originality; Research method; Main findings; Implications for theory and practice).
% Conteudo: o ABSTRACT DO ORIGINAL (Paper Training - 15 Nov 2025), redistribuido nos cinco rotulos
% (06/11/2026). Acrescentado do original: a media de 66,90 no Nivel 2 (secao 4.4 do original).
% Palavras-chave: o original nao tinha; 3 a 5, sem repetir termos do titulo (Kirkpatrick, framework,
% integrating, safety, sustainability, aviation, maintenance, waste, reduction, resource, efficiency).

\section*{Abstract}

\textbf{Paper aims:} In high-risk industries such as aviation maintenance, safety and sustainability are closely linked. Maintenance errors, a primary contributor to accidents, also generate resource waste through energy-intensive rework, material scrap and the loss of human capital. This study proposes and implements a Kirkpatrick-based framework that integrates safety and sustainability objectives in the technical training of a multinational aircraft engine maintenance facility.

\textbf{Originality:} The four-level model (reaction, learning, behavior, results) is operationalized under AS9100 and civil aviation regulation to evaluate training for safety and compliance and also for its effect on resource efficiency. Evaluation thus becomes a continuous improvement mechanism.

\textbf{Research method:} Twelve-month sequential mixed-methods case study combining post-training questionnaires, learning assessments, paired leader--trainee questionnaires, customer satisfaction ratings (product quality, technical documentation, turnaround time) and quality-system metrics.

\textbf{Main findings:} The implementation was associated with 28\% fewer training failures and 35\% fewer regulatory non-conformances, and with lower quality costs, implying avoided material and energy use in corrective actions, not yet quantified in physical units. Learning assessment (mean 66.90 out of 100) exposed content gaps that led to a content review.

\textbf{Implications for theory and practice:} A systematic training evaluation framework can turn technical instruction from a compliance obligation into a core strategy for cleaner production. The case offers a replicable model for integrated safety and sustainability in high-risk, resource-intensive operations.

\textbf{Keywords:} training evaluation; AS9100; cost of quality; human factors; cleaner production.
